\documentclass[11pt]{article}

% =========================================================
% Packages
% =========================================================
\usepackage[margin=1in]{geometry}
\usepackage{amsmath,amssymb,amsfonts}
\usepackage{booktabs}
\usepackage{graphicx}
\usepackage{microtype}
\usepackage{hyperref}
\usepackage{enumitem}
\usepackage{xcolor}
\usepackage{multirow}
\usepackage{array}
\usepackage{longtable}
\usepackage{algorithm}
\usepackage{algpseudocode}
\usepackage{natbib}

\hypersetup{
    colorlinks=true,
    linkcolor=blue,
    citecolor=blue,
    urlcolor=blue
}

\newcommand{\R}{\mathbb{R}}
\newcommand{\Prob}{\mathbb{P}}
\newcommand{\E}{\mathbb{E}}
\newcommand{\ind}{\mathbb{I}}
\newcommand{\sigmoid}{\operatorname{sigmoid}}
\newcommand{\softmax}{\operatorname{softmax}}
\newcommand{\BCE}{\mathcal{L}_{\mathrm{BCE}}}
\newcommand{\TSFM}{\mathrm{TSFM}}

% =========================================================
% Title
% =========================================================
\title{
\textbf{Direction-Aware Multimodal Transfer Learning\\
for Stock Forecasting with Time-Series Foundation Models}
}

\author{
Hoa Nguyen
}

\date{\today}

\begin{document}

\maketitle

% =========================================================
% Abstract
% =========================================================
\begin{abstract}
This proposal develops a multimodal forecasting framework for predicting the
direction of four-week-ahead stock-price movements using historical market
data and precomputed text embeddings. The central design principle is alignment
with the competition's hit-rate evaluation metric. Rather than treating the
problem primarily as point forecasting, we formulate it as binary directional
classification.

The proposed framework transfers information from competition-permitted
pretrained time-series foundation models, including Chronos, Moirai, and
TimesFM, through three progressively more adaptive mechanisms: forecast
stacking, frozen latent-representation extraction, and parameter-efficient
fine-tuning. The implemented architecture combines a frozen TimesFM price
representation with a Temporal Fusion Transformer (TFT) over historical
market covariates. Original, uncompressed text embeddings are mapped to a
shared width by family-specific adapters, contextualized jointly by
self-attention, and selected by market-conditioned cross-attention. A
cross-sectional Transformer can additionally model dependencies among stocks
at each prediction date.

The empirical study will compare simple tabular classifiers, neural
time-series models trained from scratch, zero-shot foundation-model forecasts,
and multimodal transfer-learning architectures under purged walk-forward
validation. The primary outcome is out-of-sample directional hit rate, with
particular attention to robustness across time, stocks, sectors, and market
regimes.
\end{abstract}

% =========================================================
\section{Motivation}
% =========================================================

Financial return forecasting is difficult because returns exhibit low
signal-to-noise ratios, non-stationarity, structural breaks, and weak temporal
persistence. Recent time-series foundation models have been pretrained on
large and diverse collections of time series and can therefore provide useful
inductive priors in data-limited forecasting settings.

However, several issues prevent the direct use of a pretrained numerical
forecast as the final solution.

First, the competition evaluates predictions using directional hit rate rather
than a standard regression loss such as mean squared error. A forecast can have
small numerical error while predicting the wrong direction. Conversely, a
forecast can have relatively large numerical error while correctly predicting
whether the future price will be above or below its current level.

Second, pretrained forecasting models are generally trained on numerical time
series and do not directly incorporate the competition-provided text
embeddings. The text vectors may contain information about firm-specific,
sector-level, and macroeconomic developments that is not observable from price
history alone.

Third, recent evidence suggests that time-series foundation models can be
strong forecasting priors in financial applications, but no single model
consistently dominates across all assets. Locally trained models can still
outperform pretrained models for particular stocks or regimes. This motivates
a hybrid approach that combines pretrained temporal representations,
competition-specific supervised learning, and multimodal information.

The main research question is therefore:

\begin{quote}
\emph{How can competition-permitted pretrained time-series models be adapted
and combined with precomputed text embeddings to maximize out-of-sample
directional hit rate?}
\end{quote}

% =========================================================
\section{Problem Formulation}
% =========================================================

Let $i \in \{1,\ldots,N\}$ denote a stock and let $t$ denote a prediction date.
Let $P_{i,t}$ be the closing price of stock $i$ at time $t$. The forecasting
horizon is $H$, corresponding approximately to four trading weeks.

The future cumulative log return is

\begin{equation}
R_{i,t}^{(H)}
=
\log\left(
\frac{P_{i,t+H}}{P_{i,t}}
\right).
\end{equation}

The directional target is

\begin{equation}
y_{i,t}
=
\ind\left\{
R_{i,t}^{(H)} > 0
\right\}.
\end{equation}

Thus, although the forecast is indexed by its origin date $t$, its realized
label is determined by the endpoint $t+H$. All inputs used to produce a
forecast at $t$ must be observable by the competition's prediction cutoff on
that date.

Let

\begin{equation}
X^{P}_{i,t-L+1:t}
\in
\R^{L \times F_P}
\end{equation}

denote the historical price-related feature sequence over a lookback window
of length $L$, and let

\begin{equation}
z^T_{i,t}
\in
\R^{D_T}
\end{equation}

denote the supplied precomputed text embedding or a derived aggregation of
the supplied text embeddings.

The primary model estimates

\begin{equation}
p_{i,t}
=
\Prob
\left(
y_{i,t}=1
\mid
X^P_{i,t-L+1:t},
z^T_{i,t}
\right).
\end{equation}

The predicted direction is

\begin{equation}
\widehat{y}_{i,t}
=
\ind\left\{
p_{i,t} > q
\right\},
\end{equation}

where $q$ is a classification threshold selected using validation data.

The principal evaluation metric is

\begin{equation}
\mathrm{HitRate}
=
\frac{1}{|\mathcal{D}_{\mathrm{test}}|}
\sum_{(i,t)\in\mathcal{D}_{\mathrm{test}}}
\ind
\left\{
\widehat{y}_{i,t}=y_{i,t}
\right\}.
\end{equation}

If the submission format requires a numerical future price rather than a
binary direction, probabilities can be converted into valid price forecasts
using

\begin{equation}
\widehat{P}_{i,t+H}
=
P_{i,t}
\exp
\left[
c(2p_{i,t}-1)
\right],
\end{equation}

where $c>0$ is a small fixed constant. This conversion preserves the predicted
direction:

\begin{equation}
\widehat{P}_{i,t+H} > P_{i,t}
\quad\Longleftrightarrow\quad
p_{i,t}>0.5.
\end{equation}

% =========================================================
\section{Notation and Symbol Definitions}
% =========================================================

Superscripts identify either a data modality or a model; they are not powers
unless stated otherwise. In particular, $P$ denotes the price modality, $T$
denotes the text modality, and a superscript $\top$ denotes matrix transpose.
A tilde denotes a transformed or fused representation, a hat denotes an
estimate or prediction, and a bar denotes an empirical average. The notation
used throughout the paper is defined below.

{\renewcommand{\arraystretch}{1.12}
\begin{longtable}{@{}p{0.25\textwidth}p{0.69\textwidth}@{}}
\caption{Core indices, targets, predictions, and datasets.}
\label{tab:notation-core}\\
\toprule
\textbf{Symbol} & \textbf{Concrete definition} \\
\midrule
\endfirsthead
\toprule
\textbf{Symbol} & \textbf{Concrete definition} \\
\midrule
\endhead
$i\in\{1,\ldots,N\}$ & Stock index; $N$ is the number of stocks in the competition universe. \\
$t$ & Forecast-origin date. Every input indexed by $t$ is restricted to information observable by the prediction cutoff on that date. \\
$H$ & Fixed forecast horizon in trading days; approximately four trading weeks (typically 20 trading days). \\
$L$ & Number of historical dates in the price lookback window. \\
$P_{i,t}$ & Observed closing price of stock $i$ on date $t$. \\
$\widehat P_{i,t+H}$ & Predicted closing price for stock $i$ at the endpoint $t+H$. \\
$R_{i,t}^{(H)}$ & Realized cumulative log return, $\log(P_{i,t+H}/P_{i,t})$, from the forecast origin to the endpoint. \\
$y_{i,t}$ & Direction label, equal to one when $R_{i,t}^{(H)}>0$ and zero otherwise; it becomes observable only at $t+H$. \\
$p_{i,t}$ & Model-estimated probability, computed at $t$, that $y_{i,t}=1$, equivalently that $P_{i,t+H}>P_{i,t}$. \\
$\widehat y_{i,t}$ & Binary direction forecast $\ind\{p_{i,t}>q\}$. \\
$q$; $q^\star$ & Classification threshold; $q^\star$ is the threshold selected exclusively from validation predictions. \\
$c$ & Positive log-return magnitude used only to convert a probability into a valid numerical price forecast. \\
$\mathcal D_{\mathrm{train}}$, $\mathcal D_{\mathrm{validation}}$, $\mathcal D_{\mathrm{test}}$ & Sets of stock--date observations assigned to training, validation, and testing. \\
$|\mathcal D|$ & Number of observations in dataset or index set $\mathcal D$. \\
$\ind\{\cdot\}$ & Indicator function: one if its condition is true and zero otherwise. \\
$\Prob(\cdot)$, $\E[\cdot]$, $\R^d$ & Probability, expectation, and the $d$-dimensional real vector space. \\
\bottomrule
\end{longtable}

\begin{longtable}{@{}p{0.25\textwidth}p{0.69\textwidth}@{}}
\caption{Input-feature and representation notation.}
\label{tab:notation-inputs}\\
\toprule
\textbf{Symbol} & \textbf{Concrete definition} \\
\midrule
\endfirsthead
\toprule
\textbf{Symbol} & \textbf{Concrete definition} \\
\midrule
\endhead
$X^P_{i,t-L+1:t}\in\R^{L\times F_P}$ & Ordered price-feature sequence for stock $i$ ending at $t$; rows are dates and the $F_P$ columns are price, return, volatility, volume, market, or related variables. \\
$x^P_{i,t}\in\R^{F_{\mathrm{tab}}}$ & Fixed-length tabular price vector evaluated at $t$, formed from the engineered price statistics listed in the Price Features subsection. \\
$x^C_{i,t}\in\R^{F_C}$ & Complete engineered covariate vector available at forecast origin $t$: market and volume features, calendar and Fourier features, and text-availability counts or indicators. It enters the implemented TFT as static context together with the frozen TimesFM representation. \\
$\bar x^C_{i,t}\in\R^{F_C}$ & Fold-standardized covariate vector. Missing entries of $x^C_{i,t}$ are replaced by training-fold medians and each column is centered and scaled using training-fold statistics only. \\
$X^M_{i,t-T+1:t}\in\R^{T\times V}$ & Fold-standardized historical market window supplied to the encoder-only TFT. The implementation uses $T=32$ and $V=16$ genuinely historical price, return, volatility, volume, and market-relative variables. \\
$F_P$; $F_{\mathrm{tab}}$ & Number of features per sequential time step and number of engineered features in the tabular price vector, respectively. \\
$r_{i,t}^{(k)}$ & Log return from date $t-k$ to date $t$, $\log(P_{i,t}/P_{i,t-k})$; $r_{i,t}^{(1)}$ is the return from $t-1$ to $t$. \\
$\bar r_{i,t}^{(k)}$; $\sigma_{i,t}^{(k)}$ & Mean and sample standard deviation of the $k$ one-day returns from $t-k+1$ through $t$. \\
$e^{(k)}_{i,t,j}\in\R^{D_k}$ & Supplied embedding from text-embedding family $k$ for text item $j$ associated with stock $i$ and date $t$; $D_k$ is its original dimension. \\
$n_{i,t}$ & Number of distinct text items associated with stock $i$ and date $t$. \\
$\mu_k$; $U_k\in\R^{D_k\times d_k}$ & Training-set mean and PCA loading matrix for embedding family $k$. \\
$\widetilde e^{(k)}_{i,t,j}\in\R^{d_k}$ & PCA-compressed embedding of text item $j$ from family $k$. \\
$\bar{\widetilde e}^{(k)}_{i,t}\in\R^{d_k}$ & Family-$k$ stock-date vector obtained by averaging its PCA-compressed vectors over the distinct associated text items. \\
$K_T$ & Number of supplied text-embedding families used by the model. \\
$z^T_{i,t}\in\R^{D_T}$ & Uncompressed supplied or aggregated text-feature vector available for stock $i$ by date $t$; $D_T$ is its dimension. \\
$s^T_{i,t}\in\R^{F_T}$ & Optional text-summary vector containing features such as item count, embedding norm, temporal averages, or similarity statistics. \\
$\widetilde z^T_{i,t}\in\R^{\widetilde D_T}$ & Tabular text vector obtained by concatenating the article-averaged vector from every embedding family and $s^T_{i,t}$; a missing family vector is zero. \\
$h^{T,k,a}_{i,t}$; $h^T_{i,t}$; $H^T_{i,t}$ & Adapted token for embedding family $k$ and text field $a$, market-conditioned pooled text vector, and the sequence of all available family--field tokens, respectively. \\
$A_T$; $J$; $d_T$ & Number of text fields, maximum family--field tokens $J=K_TA_T$, and shared text width. The implementation uses $d_T=384$. \\
\bottomrule
\end{longtable}

\begin{longtable}{@{}p{0.25\textwidth}p{0.69\textwidth}@{}}
\caption{Model, fusion, and optimization notation.}
\label{tab:notation-models}\\
\toprule
\textbf{Symbol} & \textbf{Concrete definition} \\
\midrule
\endfirsthead
\toprule
\textbf{Symbol} & \textbf{Concrete definition} \\
\midrule
\endhead
$f_\theta$, $g_\psi$ & Predictive mapping and task-specific head with learned parameter collections $\theta$ and $\psi$. Subscripts identify the corresponding model. \\
$m$; $M_{\mathrm{mod}}$ & Model index and number of models or ensemble members. \\
$\theta_m^0$; $\theta_0$ & Frozen pretrained parameters of foundation model $m$, or of the selected backbone when no model index is needed. \\
$H^{P,m}_{i,t}\in\R^{M_m\times d_m}$ & Hidden price-token matrix from foundation model $m$; $M_m$ and $d_m$ are its token count and hidden width. \\
$H^P_{i,t}\in\R^{M_P\times d}$; $h^P_{i,t}$ & Price-token matrix in the shared latent space and its pooled vector; $M_P$ is the number of price tokens. \\
$h^C_{i,t}$ & Joint market state produced by the encoder-only TFT from the frozen TimesFM vector, static origin covariates, and the historical market window. \\
$p^{(m)}_{i,t}$; $a^{(m)}_{i,t}$ & Model-$m$ probability of an upward endpoint and its signed point-forecast log-return score. \\
$S_m$; $\widehat P^{(m,s)}_{i,t+H}$ & Number of forecast samples from model $m$ and its $s$th sampled endpoint price. \\
$u^{(m)}_{i,t}$; $h^{\mathrm{FM}}_{i,t}$ & Model-$m$ uncertainty summaries and the concatenated foundation-model forecast-feature vector. \\
$W_0$, $W$, $\Delta W$ & Original pretrained weight matrix, adapted weight matrix, and learned LoRA update. \\
$A_{\mathrm{LoRA}}$, $B_{\mathrm{LoRA}}$, $r$ & LoRA down-projection, up-projection, and rank, with $r\ll\min(d_{\mathrm{in}},d_{\mathrm{out}})$. \\
$d_{\mathrm{in}}$; $d_{\mathrm{out}}$ & Input and output widths of the adapted weight matrix. \\
$\phi$; $\phi_{\mathrm{LoRA}}$ & Optional adapter parameters and the collection of all trainable LoRA parameters. \\
$W_T,b_T,\phi_T$ & Family-specific raw-text adapter parameters and nonlinear projection mapping. \\
$q^C,K^T,V^T$ & Single market-state query and contextual text keys/values used in the implemented market-to-text cross-attention. \\
$W_Q,W_K,W_V$; $d_{\mathrm{att}}$ & Learned attention projections and per-head query/key width. \\
$\widetilde H^T_{i,t}$; $\widetilde H^P_{i,t}$ & Cross-attended text-token and price-token matrices. \\
$g^P_{i,t},g^T_{i,t}$ & Elementwise price and text gates in $[0,1]^d$. \\
$W_g,b_g,U_g,c_g$ & Trainable matrices and biases used to construct the two gates. \\
$\widetilde h^P_{i,t}$; $\widetilde h^T_{i,t}$ & Price and text vectors after elementwise gating. \\
$c^{(\ell)}_{i,\tau}$; $B_{\mathrm{TCN}}$ & TCN hidden vector at date $\tau$ after residual block $\ell$, and the total number of TCN blocks. \\
$K_{\mathrm{TCN}}$; $\delta_\ell$ & TCN convolution-kernel width and dilation of block $\ell$, with $\delta_\ell=2^{\ell-1}$. \\
$h^{\mathrm{TCN}}_{i,t}$; $u^{\mathrm{TCN}}_{i,t}$ & Final causal TCN price state at forecast origin $t$, and its fused price--text representation. \\
$p^{\mathrm{TCN}}_{i,t}$ & Probability of an upward endpoint produced by the train-from-scratch TCN with text fusion. \\
$v_{i,\tau}$; $h^{\mathrm{GRU}}_{i,\tau}$ & Multimodal input vector and recurrent GRU hidden state at historical date $\tau$. \\
$a^{\mathrm{upd}}_{i,\tau}$; $a^{\mathrm{rst}}_{i,\tau}$ & GRU update and reset gates. \\
$\bar h^{\mathrm{GRU}}_{i,\tau}$; $d_g$ & Candidate GRU state and GRU hidden-state width. \\
$\odot$; $[a;b]$ & Elementwise multiplication and vector concatenation. \\
$u_{i,t}\in\R^{d_u}$ & Fused representation of stock $i$ at date $t$; $d_u$ is its dimension. \\
$U_t\in\R^{N\times d_u}$; $\widetilde U_t$ & Matrix of all stock representations at date $t$ and its cross-sectionally contextualized version. \\
$\widetilde u_{i,t}$; $w,b$ & Row of $\widetilde U_t$ for stock $i$, and the final classifier's weight vector and scalar bias. \\
$X\in\R^{L\times F}$; $x^{(f)}$ & Generic iTransformer input with $F$ variables, and the length-$L$ trajectory of variable $f$. \\
$W_{\mathrm{emb}},b_{\mathrm{emb}},h_f^{(0)}$ & iTransformer trajectory-embedding matrix, bias, and initial token for variable $f$. \\
$\mathcal A\in\R^{F\times F}$ & iTransformer attention-weight matrix across variables. \\
\bottomrule
\end{longtable}

\begin{longtable}{@{}p{0.25\textwidth}p{0.69\textwidth}@{}}
\caption{Loss, validation, ensemble, and evaluation notation.}
\label{tab:notation-evaluation}\\
\toprule
\textbf{Symbol} & \textbf{Concrete definition} \\
\midrule
\endfirsthead
\toprule
\textbf{Symbol} & \textbf{Concrete definition} \\
\midrule
\endhead
$\mathcal B$; $B=|\mathcal B|$ & Training mini-batch of stock--date observations and its size. \\
$q_{i,t}$; $p^{\mathrm{LR}}_{i,t}$ & Logistic-regression input vector and its predicted probability of an upward endpoint. \\
$\beta_0$; $\beta$ & Unpenalized logistic-regression intercept and coefficient vector. \\
$\lambda_2$; $\mathcal L_{\mathrm{LR}}$ & Nonnegative $L_2$ regularization strength and the penalized logistic-regression objective. \\
$\BCE$ & Mean binary cross-entropy over the current mini-batch. \\
$\mathcal L_{\mathrm{weighted}}$; $\mathcal L$ & Confidence-weighted classification loss and the combined classification--regression objective. \\
$w_{i,t}$; $w_{\max}$ & Confidence weight assigned to observation $(i,t)$ and its upper cap. \\
$\alpha$; $\widehat\sigma_{i,t}^{(H)}$; $\epsilon$ & Weighting strength, estimated scale of the $H$-day return, and a small positive numerical stabilizer. \\
$\widehat R_{i,t}^{(H)}$; $g_R$ & Predicted $H$-day log return and its auxiliary regression head. \\
$\lambda_R$; $\lambda_D$ & Nonnegative weights on the auxiliary return loss and distillation loss. \\
$\operatorname{Huber}$; $\operatorname{KL}$ & Huber regression loss and Kullback--Leibler divergence. \\
$p^{\mathrm{teacher}}_{i,t}$; $p^{\mathrm{student}}_{i,t}$ & Ensemble teacher probability and student-model probability used for distillation. \\
$w_S,b_S$ & Weight vector and scalar bias of the standalone student classifier. \\
$\mathcal L_{\mathrm{distill}}$ & Combined binary-classification and teacher--student distillation objective. \\
$\alpha_m$ & Nonnegative weight for model $m$; the relevant model weights sum to one. \\
$\kappa$; $K_{\mathrm{fold}}$ & Walk-forward fold index and total number of validation folds. \\
$T_\kappa$; $V_\kappa$; $G$ & Last training date, last validation date, and purge/embargo length for fold $\kappa$, with $G\geq H$. \\
$S_{\mathrm{robust}}$; $\gamma$ & Mean fold hit rate minus $\gamma$ times its fold-to-fold standard deviation, and the nonnegative instability penalty. \\
$\mathrm{HitRate}_\kappa(q)$ & Hit rate in validation fold $\kappa$ when probabilities are converted to labels using threshold $q$. \\
$\mathcal Q$ & Finite candidate set searched when selecting $q^\star$. \\
$p^{\mathrm{ens}}_{i,t}$; $D_{i,t}$ & Ensemble probability and standard deviation of member probabilities (model disagreement). \\
$\widehat h$; $n$ & Observed hit rate and number of evaluated predictions used in its elementary standard-error calculation. \\
$A,B$ (model labels) & Two prediction methods being compared in the paired statistical analysis; these labels are unrelated to the LoRA matrices. \\
$N_t$; $d_t$ & Number of evaluated stocks on date $t$ and the date-level difference in hit rate between models $A$ and $B$. \\
\bottomrule
\end{longtable}
}

% =========================================================
\section{Competition Constraints}
% =========================================================

The proposed methodology is designed around the following restrictions:

\begin{enumerate}[label=(\arabic*)]
    \item No external datasets may be used.
    \item External pretrained models are prohibited except for the explicitly
    permitted models Chronos, Moirai, TimesFM, and TabPFN.
    \item Text may only be used through the precomputed embedding vectors
    supplied by the competition.
    \item Raw text may not be processed through prompts, external APIs, or
    additional language models.
    \item All trainable preprocessing procedures, including normalization,
    feature selection, and dimensionality reduction, must be fitted using only
    the competition training data.
\end{enumerate}

Before implementing model fine-tuning, the rules must be checked to determine
whether permission to use the four pretrained models includes:

\begin{itemize}
    \item zero-shot inference only;
    \item extraction of internal hidden representations;
    \item partial or full fine-tuning;
    \item parameter-efficient adaptation such as LoRA.
\end{itemize}

The proposed experiment therefore separates methods by their level of
adaptation. If fine-tuning is disallowed, the frozen-forecast and frozen-feature
variants remain valid candidate methods.

% =========================================================
\section{Input Construction}
% =========================================================

\subsection{Price Features}

For each stock and date, the raw price sequence will be transformed into
stationary or approximately stationary variables. Candidate features include:

\begin{align}
r_{i,t}^{(1)}
&=
\log\left(\frac{P_{i,t}}{P_{i,t-1}}\right),\\
r_{i,t}^{(k)}
&=
\log\left(\frac{P_{i,t}}{P_{i,t-k}}\right),
\qquad
k\in\{5,10,20,40,60\},\\
\bar r_{i,t}^{(k)}
&=
\frac{1}{k}
\sum_{j=0}^{k-1}r_{i,t-j}^{(1)},\\
\sigma_{i,t}^{(k)}
&=
\sqrt{
\frac{1}{k-1}
\sum_{j=0}^{k-1}
\left(
r_{i,t-j}^{(1)}
-
\bar r_{i,t}^{(k)}
\right)^2
}.
\end{align}

Additional candidate features include:

\begin{itemize}
    \item rolling high--low range;
    \item close-to-open return;
    \item rolling drawdown;
    \item moving-average gaps;
    \item momentum acceleration;
    \item rolling return skewness;
    \item rolling return kurtosis;
    \item volume level and volume changes, when supplied;
    \item stock return relative to the supplied stock universe;
    \item cross-sectional return rank;
    \item sector-relative return, if sector labels are supplied.
\end{itemize}

All cross-sectional features will be constructed solely from the competition
universe.

For a tabular classifier, these quantities are evaluated at the forecast
origin and collected into the fixed-length vector

\begin{equation}
x^P_{i,t}
=
\operatorname{Tabularize}
\left(X^P_{i,t-L+1:t}\right)
\in\R^{F_{\mathrm{tab}}},
\end{equation}

where $\operatorname{Tabularize}$ denotes the deterministic construction of
the listed lagged returns, rolling statistics, market-relative variables, and
other engineered price features. Its parameters are fitted using training
data only.

\subsection{Text Embeddings}

Suppose the competition provides one or more allowed embedding vectors for
stock $i$ and date $t$. For embedding family $k$, let

\begin{equation}
e^{(k)}_{i,t,j}
\in
\R^{D_k}
\end{equation}

denote the embedding of text item $j$.

For each embedding family, a dimensionality-reduction mapping is fitted using
training observations only and is applied separately to every text item:

\begin{equation}
\widetilde e^{(k)}_{i,t,j}
=
U_k^\top
\left(
e^{(k)}_{i,t,j}
-
\mu_k
\right),
\end{equation}

where $U_k$ may be obtained from principal component analysis, with candidate
retained dimensions

\begin{equation}
d_k \in \{16,32,64,128\}.
\end{equation}

For a tabular model, the compressed articles are averaged only within the same
embedding family:

\begin{equation}
\bar{\widetilde e}^{(k)}_{i,t}
=
\frac{1}{n_{i,t}}
\sum_{j=1}^{n_{i,t}}
\widetilde e^{(k)}_{i,t,j}
\in\R^{d_k}.
\end{equation}

Embedding families are never averaged with one another because independently
fitted PCA axes are not aligned across families.  If no usable article is
available for family $k$, its aggregated vector is set to zero, while separate
count and availability features identify the missing-text case.

For neural models, dimensionality reduction may instead be learned through a
trainable projector:

\begin{equation}
h^{T,k}_{i,t,j}
=
\rho
\left(
W_k e^{(k)}_{i,t,j}+b_k
\right),
\end{equation}

where $W_k$ and $b_k$ are the family-specific projection parameters and
$\rho$ is a fixed nonlinear activation.

For tabular models, the family-specific article means and optional summary
statistics are concatenated as

\begin{equation}
\widetilde z^T_{i,t}
=
\left[
\bar{\widetilde e}^{(1)}_{i,t};
\ldots;
\bar{\widetilde e}^{(K_T)}_{i,t};
s^T_{i,t}
\right]
\in\R^{\widetilde D_T},
\end{equation}

where $s^T_{i,t}$ contains text counts and availability indicators and
$\widetilde D_T=\sum_{k=1}^{K_T}d_k+F_T$ is the dimension of the resulting
concatenated vector.  Attention-based models instead retain the item-level
vectors $\widetilde e^{(k)}_{i,t,j}$ as separate tokens.

% =========================================================
\section{Baseline Models}
% =========================================================

\subsection{Naive Directional Baselines}

The following baselines are necessary to determine whether the learned models
provide genuine incremental directional information:

\begin{enumerate}
    \item always predict an upward movement;
    \item predict the sign of the previous one-day return;
    \item predict the sign of the previous 20-day return;
    \item predict the sign of the stock's excess return relative to the market;
    \item sector momentum, if sector information is available.
\end{enumerate}

The always-up baseline is particularly important because long-run positive
equity drift may result in an unconditional upward frequency above 50\%.

\subsection{Tabular Classification}

LightGBM or CatBoost will be trained on engineered price variables and
compressed text embeddings:

\begin{equation}
p^{\mathrm{tab}}_{i,t}
=
f_{\theta_{\mathrm{tab}}}
\left(
x^P_{i,t},
\widetilde z^T_{i,t}
\right)
=
\widehat{\Prob}_{\theta_{\mathrm{tab}}}
\left(
P_{i,t+H}>P_{i,t}
\mid
x^P_{i,t},\widetilde z^T_{i,t}
\right).
\end{equation}

Here $p^{\mathrm{tab}}_{i,t}$ is indexed by $t$ because it is produced at the
forecast origin; the event it predicts is explicitly evaluated at $t+H$.

This model provides a strong low-variance benchmark and can capture nonlinear
interactions between price conditions and embedding dimensions.

TabPFN will also be tested when computationally feasible:

\begin{equation}
p^{\mathrm{TabPFN}}_{i,t}
=
f_{\mathrm{TabPFN}}
\left(
x^P_{i,t},
\widetilde z^T_{i,t};
\mathcal{D}_{\mathrm{train}}
\right).
\end{equation}

\subsection{Train-from-Scratch Sequential Baselines}

The principal train-from-scratch models will include:

\begin{itemize}
    \item temporal convolutional network;
    \item gated recurrent unit;
    \item PatchTST-style encoder;
    \item iTransformer.
\end{itemize}

Each model will be equipped with a binary classification head rather than only
a numerical forecasting head.

% =========================================================
\section{Pretrained Time-Series Foundation Models}
% =========================================================

Let $m$ index a permitted pretrained model:

\begin{equation}
m
\in
\{
\mathrm{Chronos},
\mathrm{Moirai},
\mathrm{TimesFM}
\}.
\end{equation}

Let $\theta_m^0$ denote its pretrained parameters.

\subsection{Frozen Forecast Features}

The simplest transfer strategy uses each model without updating its parameters.
For model $m$, obtain either a point forecast, forecast samples, or forecast
quantiles.

If forecast samples are available, define

\begin{equation}
p^{(m)}_{i,t}
=
\frac{1}{S_m}
\sum_{s=1}^{S_m}
\ind
\left\{
\widehat P^{(m,s)}_{i,t+H}
>
P_{i,t}
\right\}.
\end{equation}

Here $S_m$ is the number of predictive samples drawn from model $m$, and
$\widehat P^{(m,s)}_{i,t+H}$ is sample $s$ of its endpoint-price forecast.

This quantity estimates the model-implied probability of a positive terminal
return.

If only a point forecast is available, define the signed score

\begin{equation}
a^{(m)}_{i,t}
=
\log
\left(
\frac{
\widehat P^{(m)}_{i,t+H}
}{
P_{i,t}
}
\right).
\end{equation}

Foundation-model features are collected as

\begin{equation}
h^{\mathrm{FM}}_{i,t}
=
\left[
p^{(\mathrm{Chronos})}_{i,t},
p^{(\mathrm{Moirai})}_{i,t},
p^{(\mathrm{TimesFM})}_{i,t},
u^{(\mathrm{Chronos})}_{i,t},
u^{(\mathrm{Moirai})}_{i,t},
u^{(\mathrm{TimesFM})}_{i,t}
\right],
\end{equation}

where $u^{(m)}_{i,t}$ represents uncertainty statistics such as forecast
dispersion or interquantile range.

These features are supplied to a competition-trained classifier:

\begin{equation}
p_{i,t}
=
f_\psi
\left(
x^P_{i,t},
z^T_{i,t},
h^{\mathrm{FM}}_{i,t}
\right).
\end{equation}

\subsection{Frozen Latent Representation Transfer}

A richer use of pretraining is to extract internal hidden states rather than
only the final forecast.

For model $m$, let

\begin{equation}
H^{P,m}_{i,t}
=
f^{\mathrm{hidden}}_{\theta_m^0}
\left(
X^P_{i,t-L+1:t}
\right)
\in
\R^{M_m \times d_m}
\end{equation}

denote its sequence of hidden price tokens.

The pooling operator is defined as mean pooling over the token dimension:

\begin{equation}
\operatorname{Pool}
\left(
H^{P,m}_{i,t}
\right)
=
\frac{1}{M_m}
\sum_{\ell=1}^{M_m}
H^{P,m}_{i,t}[\ell,:]
\in
\R^{d_m},
\end{equation}

where $H^{P,m}_{i,t}[\ell,:]$ is the $d_m$-dimensional hidden vector of
token $\ell$. If padding tokens are present, the average is taken only over
valid, non-padding tokens and $M_m$ denotes their count.

The pretrained backbone remains frozen, while a competition-specific pooling
and classification network is trained:

\begin{align}
h^{P,m}_{i,t}
&=
\operatorname{Pool}
\left(
H^{P,m}_{i,t}
\right),\\
p_{i,t}
&=
\sigmoid
\left[
g_\psi
\left(
h^{P,m}_{i,t},
\bar x^C_{i,t},
z^T_{i,t}
\right)
\right].
\end{align}

This approach tests whether the foundation model's latent representation
contains useful information that is discarded by its original numerical
forecast head.

\subsection{Parameter-Efficient Adaptation}

If permitted by the competition rules, the main adaptive model will use
low-rank adaptation. For a pretrained weight matrix $W_0$, LoRA defines

\begin{equation}
W
=
W_0 + \Delta W,
\qquad
\Delta W = B_{\mathrm{LoRA}}A_{\mathrm{LoRA}},
\end{equation}

where

\begin{equation}
A_{\mathrm{LoRA}}\in\R^{r\times d_{\mathrm{in}}},
\qquad
B_{\mathrm{LoRA}}\in\R^{d_{\mathrm{out}}\times r},
\end{equation}

and $r$ is much smaller than the original matrix dimensions.

The backbone becomes

\begin{equation}
H^P_{i,t}
=
f_{\theta_0,\phi_{\mathrm{LoRA}}}
\left(
X^P_{i,t-L+1:t}
\right),
\end{equation}

where $\theta_0$ remains frozen and only $\phi_{\mathrm{LoRA}}$ is updated.

This allows the pretrained model to adapt to:

\begin{itemize}
    \item financial return dynamics;
    \item the four-week prediction horizon;
    \item the directional classification objective;
    \item the specific stock universe;
    \item interactions with text-derived information.
\end{itemize}

% =========================================================
\section{Proposed Multimodal Architecture}
% =========================================================

\subsection{Pretrained Price Encoder}

The preferred initial backbone is TimesFM because it provides a pretrained
temporal representation and can potentially be adapted using a limited number
of trainable parameters.

The price representation is

\begin{equation}
H^P_{i,t}
=
f_{\theta_0,\phi}
\left(
X^P_{i,t-L+1:t}
\right)
\in
\R^{M_P\times d},
\end{equation}

where:

\begin{itemize}
    \item $\theta_0$ denotes frozen pretrained parameters;
    \item $\phi$ denotes optional LoRA or adapter parameters;
    \item $M_P$ denotes the number of temporal or patch tokens;
    \item $d$ denotes the shared latent width of each token.
\end{itemize}

When internal representation extraction is unavailable, a train-from-scratch
TCN or PatchTST encoder will be used instead.

\subsection{TFT Market Encoder}

The implemented model pools a frozen TimesFM hidden representation
$h^P_{i,t}$ extracted from the historical \texttt{ret\_20} series. TimesFM
uses up to 512 observations and requires at least 64 valid observations; the
result is cached and is not updated during classifier training. The pooled
state and all current-origin covariates form a static context:

\begin{equation}
c_{i,t}
=
\operatorname{GRN}_{\mathrm{static}}
\left([h^P_{i,t};\bar x^C_{i,t}]\right).
\end{equation}

In parallel, the TFT consumes the preceding 32 dates of 16 historical market
variables, $X^M_{i,t-31:t}$. Each scalar variable is encoded by its own gated
residual network (GRN). A context-conditioned variable-selection network
produces weights $\alpha_{i,\tau,v}$:

\begin{align}
\alpha_{i,\tau}
&=
\softmax\!\left(
\operatorname{GRN}_{\mathrm{weight}}
(X^M_{i,\tau},c_{i,t})
\right),\\
s_{i,\tau}
&=
\sum_{v=1}^{V}
\alpha_{i,\tau,v}
\operatorname{GRN}_{v}(X^M_{i,\tau,v}).
\end{align}

The selected sequence is passed through an LSTM while retaining the hidden
state at every time step:

\begin{equation}
L_{i,t}
=
\operatorname{LSTM}(s_{i,t-31:t})
\in\R^{T\times d_C}.
\end{equation}

The complete LSTM output is transformed by a gated linear unit (GLU) and
combined with the selected-variable sequence through a normalized residual
connection:

\begin{equation}
R_{i,t}
=
\operatorname{LayerNorm}
\left(
s_{i,t-31:t}
+
\operatorname{GLU}(W_L L_{i,t}+b_L)
\right)
\in\R^{T\times d_C}.
\end{equation}

Multi-head self-attention then processes the complete recurrent sequence.
Left-padded positions are excluded from the attention keys and values by the
padding mask:

\begin{align}
Z_{i,t}
&=
\operatorname{MHA}
\left(
Q=R_{i,t},
K=R_{i,t},
V=R_{i,t};
m^M_{i,t}
\right),\\
A_{i,t}
&=
\operatorname{LayerNorm}
\left(
R_{i,t}
+
\operatorname{GLU}(W_A Z_{i,t}+b_A)
\right)
\in\R^{T\times d_C}.
\end{align}

Because the windows are left padded and end at the forecast origin, the code
selects the final \emph{post-attention} state $A_{i,t,-1}$, rather than the raw
final LSTM state. The static context then conditions the output GRN:

\begin{equation}
h^C_{i,t}
=
\operatorname{GRN}_{\mathrm{out}}
\left(
A_{i,t,-1},
c_{i,t}
\right)
\in\R^{d_C}.
\end{equation}

Thus the TFT is the trainable market encoder, while TimesFM supplies a frozen
price prior within its static context. No unavailable future market variable
is passed to the TFT.

\subsection{Unified Raw-Embedding Text Encoder}

The text branch consumes the original supplied embeddings directly. It does
not use PCA, mean-pool all articles into one vector, or train a separate
target-supervised TextCoder. For family $k$ and text field $a$, the raw vector
$e^{(k)}_{i,t,a}\in\R^{D_k}$ is mapped to the common width $d_T=384$:

\begin{equation}
h^{T,k,a}_{i,t}
=
\operatorname{Dropout}\!\left(
\operatorname{GELU}\!\left(
W_k\operatorname{LayerNorm}(e^{(k)}_{i,t,a})+b_k
\right)
\right)
+\rho_k+\eta_a,
\end{equation}

where $\rho_k$ and $\eta_a$ are learned family and field identity embeddings.
The family-specific adapters allow heterogeneous original dimensions while
preserving one token for every available family--field pair. The tokens are
concatenated along the sequence dimension and contextualized jointly:

\begin{equation}
\widehat H^T_{i,t}
=
\operatorname{TransformerEncoder}
\left(
[h^{T,k,a}_{i,t}]_{k,a};
m^T_{i,t}
\right)
\in\R^{J\times d_T},
\end{equation}

where $m^T_{i,t}$ masks unavailable pairs and $J=K_TA_T$. The implementation
uses one text self-attention layer with four heads by default.

\subsection{Market-Conditioned Cross-Attention and Fusion}

The TFT market state supplies one query, while the contextualized raw-text
tokens supply keys and values:

\begin{align}
q^C_{i,t}
&=
W_Qh^C_{i,t}\in\R^{1\times d_T},\\
\widehat h^T_{i,t}
&=
\operatorname{MHA}
\left(
Q=q^C_{i,t},
K=\widehat H^T_{i,t},
V=\widehat H^T_{i,t};
m^T_{i,t}
\right)
\in\R^{1\times d_T},\\
h^T_{i,t}
&=
\operatorname{GRN}_{T}
\left(
\operatorname{LayerNorm}
(q^C_{i,t}+\widehat h^T_{i,t}),
h^C_{i,t}
\right).
\end{align}

The attention weights therefore answer which embedding families and text
fields are relevant under the current market state. Masked tokens receive no
weight. The active implementation requires at least one available text token
for every modeled row and does not insert a learned null token.

The market and selected text states are concatenated, projected with GELU and
dropout, and processed by residual MLP blocks:

\begin{align}
u^{(0)}_{i,t}
&=
\operatorname{Dropout}\!\left(
\operatorname{GELU}\!\left(
W_F[h^C_{i,t};h^T_{i,t}]+b_F
\right)
\right),\\
u^{(\ell)}_{i,t}
&=
\operatorname{ResidualMLP}_{\ell}(u^{(\ell-1)}_{i,t}),\\
p_{i,t}
&=
\sigmoid\!\left(
w_o^\top\operatorname{LayerNorm}(u^{(B_F)}_{i,t})+b_o
\right).
\end{align}

TimesFM remains frozen. The TFT, family adapters, identity embeddings, text
self-attention, cross-attention, fusion blocks, and classifier are optimized
jointly within each expanding training fold. Imputation and scaling parameters
are also fitted within the fold. Hyperparameters and the decision threshold
are selected from out-of-fold predictions before the chosen architecture is
refitted on all labeled training observations.

\subsection{TCN (Temporal Convolution Network) with Text Fusion}

The TCN with text fusion is a supervised, train-from-scratch model rather than
a pretrained model.  It uses causal dilated convolutions to encode the price
history and combines the resulting price state with the supplied text
embedding.  First project each historical price-feature vector into the TCN
hidden width:

\begin{equation}
c^{(0)}_{i,\tau}
=
W_{\mathrm{in}}x^P_{i,\tau}+b_{\mathrm{in}},
\qquad
\tau\in\{t-L+1,\ldots,t\}.
\end{equation}

For residual block $\ell\in\{1,\ldots,B_{\mathrm{TCN}}\}$, use the causal
dilated update

\begin{equation}
c^{(\ell)}_{i,\tau}
=
c^{(\ell-1)}_{i,\tau}
+
\rho\left(
\sum_{q=0}^{K_{\mathrm{TCN}}-1}
W^{(\ell,q)}_{\mathrm{TCN}}
c^{(\ell-1)}_{i,\tau-q\delta_\ell}
+b^{(\ell)}_{\mathrm{TCN}}
\right),
\qquad
\delta_\ell=2^{\ell-1},
\end{equation}

where $B_{\mathrm{TCN}}$ is the number of residual blocks,
$K_{\mathrm{TCN}}$ is the convolution-kernel width, $\delta_\ell$ is the
dilation, and values before $t-L+1$ are zero padded.  The activation $\rho$
may be ReLU or GELU.  Because the convolution is causal, the representation at
$t$ uses no observations after the forecast origin.  The final price state is

\begin{equation}
h^{\mathrm{TCN}}_{i,t}
=
c^{(B_{\mathrm{TCN}})}_{i,t}.
\end{equation}

The basic concatenation-based text fusion model is

\begin{align}
u^{\mathrm{TCN}}_{i,t}
&=
\rho\left(
W_{\mathrm{fuse}}
\left[
h^{\mathrm{TCN}}_{i,t};h^T_{i,t}
\right]
+b_{\mathrm{fuse}}
\right),\\
p^{\mathrm{TCN}}_{i,t}
&=
\sigmoid\left(
w_{\mathrm{TCN}}^\top u^{\mathrm{TCN}}_{i,t}
+b_{\mathrm{TCN}}
\right).
\label{simple_concatination}
\end{align}

Here $p^{\mathrm{TCN}}_{i,t}$ estimates
$\Pr(P_{i,t+H}>P_{i,t})$.  All TCN, text-projection, fusion, and classification
parameters are learned only from the competition training data.  In the gated
variant, the concatenation projector is replaced by the elementwise gates
defined below.  In the cross-attention variant, the sequence
$[c^{(B_{\mathrm{TCN}})}_{i,t-L+1},\ldots,
c^{(B_{\mathrm{TCN}})}_{i,t}]$ serves as the price tokens attended to by the
text representation.

\subsection{Encoder--Fusion Summary}

The active Stage~2 architecture is the following hybrid:

\begin{equation}
\begin{aligned}
h^P_{i,t}
&=
\operatorname{Pool}\!\left(
\operatorname{TimesFM}_{\theta_0}
(X^{\mathrm{ret20}}_{i,t-L+1:t})
\right),\\
h^C_{i,t}
&=
\operatorname{TFT}
\left([h^P_{i,t};\bar x^C_{i,t}],X^M_{i,t-31:t}\right),\\
\widehat H^T_{i,t}
&=
\operatorname{TextSelfAttention}
\left([\operatorname{Adapter}_k(e^{(k)}_{i,t,a})
+\rho_k+\eta_a]_{k,a}\right),\\
h^T_{i,t}
&=
\operatorname{CrossAttention}
\left(Q=W_Qh^C_{i,t},K=\widehat H^T_{i,t},
V=\widehat H^T_{i,t}\right),\\
u_{i,t}
&=
\operatorname{ResidualFusion}([h^C_{i,t};h^T_{i,t}]).
\end{aligned}
\end{equation}

TimesFM is pretrained and frozen. The TFT is not a replacement for TimesFM:
it learns from the 32-step historical market window while using the frozen
price state and current-origin covariates as context. The TFT, raw-embedding
adapters, text self-attention, market-to-text cross-attention, residual fusion
blocks, and classifier are trained jointly on competition data.

\subsection{Why Cross-Attention is Appropriate}

Simple concatenation,

\begin{equation}
[h^P_{i,t};h^T_{i,t}],
\end{equation}

allows an MLP to learn interactions, but it assigns no explicit,
market-dependent importance to the individual family--field text tokens.
The implemented cross-attention instead projects the TFT market state into one
query and attends over the self-contextualized text-token sequence:

\begin{equation}
h^T_{i,t}
=
\operatorname{Attention}
\left(
Q=W_Qh^C_{i,t},
K=\widehat H^T_{i,t},
V=\widehat H^T_{i,t}
\right).
\end{equation}

This operation asks: given the present TimesFM--TFT market state, which
available embedding families and text fields are most relevant for the
directional prediction? It produces one vector of shape $1\times d_T$ per
example before the singleton dimension is removed. Attention performs a
weighted sum of the value vectors; it is not a multiplication of the query by
the attention weights.

\subsection{Gated Fusion Alternatives}

\subsubsection{Elementwise Cross-Modal Gating}

Because full cross-attention may overfit in a limited-data environment, a
lower-parameter gated fusion module will be evaluated:

\begin{align}
g^P_{i,t}
&=
\sigmoid
\left(
W_g h^T_{i,t}+b_g
\right),\\
\widetilde h^P_{i,t}
&=
g^P_{i,t}
\odot
h^P_{i,t},\\
g^T_{i,t}
&=
\sigmoid
\left(
U_g h^P_{i,t}+c_g
\right),\\
\widetilde h^T_{i,t}
&=
g^T_{i,t}
\odot
h^T_{i,t}.
\end{align}

The fused stock representation is

\begin{equation}
u_{i,t}
=
\left[
\widetilde h^P_{i,t};
\widetilde h^T_{i,t};
h^{\mathrm{FM}}_{i,t}
\right].
\end{equation}

This module is an elementwise cross-modal gate, not a recurrent unit.

\subsubsection{GRU-Based Temporal Fusion}

A genuine gated recurrent unit (GRU) will be evaluated as a train-from-scratch
alternative when aligned price and text features are available throughout the
lookback window. For each historical date
$\tau\in\{t-L+1,\ldots,t\}$, define the multimodal input

\begin{equation}
v_{i,\tau}
=
\left[
x^P_{i,\tau};
\widetilde z^T_{i,\tau}
\right].
\end{equation}

Starting from $h^{\mathrm{GRU}}_{i,t-L}=0\in\R^{d_g}$, the recurrent update is

\begin{align}
a^{\mathrm{upd}}_{i,\tau}
&=
\sigmoid\left(
W_{\mathrm{upd}}v_{i,\tau}
+U_{\mathrm{upd}}h^{\mathrm{GRU}}_{i,\tau-1}
+b_{\mathrm{upd}}
\right),\\
a^{\mathrm{rst}}_{i,\tau}
&=
\sigmoid\left(
W_{\mathrm{rst}}v_{i,\tau}
+U_{\mathrm{rst}}h^{\mathrm{GRU}}_{i,\tau-1}
+b_{\mathrm{rst}}
\right),\\
\bar h^{\mathrm{GRU}}_{i,\tau}
&=
\tanh\left(
W_hv_{i,\tau}
+U_h\left(
a^{\mathrm{rst}}_{i,\tau}
\odot
h^{\mathrm{GRU}}_{i,\tau-1}
\right)
+b_h
\right),\\
h^{\mathrm{GRU}}_{i,\tau}
&=
\left(1-a^{\mathrm{upd}}_{i,\tau}\right)
\odot h^{\mathrm{GRU}}_{i,\tau-1}
+a^{\mathrm{upd}}_{i,\tau}
\odot\bar h^{\mathrm{GRU}}_{i,\tau}.
\end{align}

Here $a^{\mathrm{upd}}_{i,\tau}$ and $a^{\mathrm{rst}}_{i,\tau}$ are the
update and reset gates, $\bar h^{\mathrm{GRU}}_{i,\tau}$ is the candidate
state, $d_g$ is the hidden-state width, and all $W$, $U$, and $b$ terms in
these equations are trainable GRU parameters. The final recurrent state is
combined with the foundation-model features as

\begin{align}
u^{\mathrm{GRU}}_{i,t}
&=
\left[
h^{\mathrm{GRU}}_{i,t};
h^{\mathrm{FM}}_{i,t}
\right],\\
p^{\mathrm{GRU}}_{i,t}
&=
\sigmoid\left(
w_{\mathrm{GRU}}^\top u^{\mathrm{GRU}}_{i,t}
+b_{\mathrm{GRU}}
\right),
\end{align}

where $u^{\mathrm{GRU}}_{i,t}$ is the fused GRU representation and
$w_{\mathrm{GRU}}$ and $b_{\mathrm{GRU}}$ are the classification-head
parameters. Unlike elementwise gating, this alternative explicitly models
the temporal order of the multimodal inputs.

% =========================================================
\section{Cross-Sectional Stock Interaction}
% =========================================================

At each date $t$, collect the fused representations of all $N$ stocks:

\begin{equation}
U_t
=
\left[
u_{1,t},
u_{2,t},
\ldots,
u_{N,t}
\right]^\top
\in
\R^{N\times d_u}.
\end{equation}

A Transformer encoder can model dependencies among stocks:

\begin{equation}
\widetilde U_t
=
\operatorname{TransformerEncoder}
\left(
U_t
\right).
\end{equation}

For stock $i$, the probability of an upward movement is

\begin{equation}
p_{i,t}
=
\sigmoid
\left(
w^\top
\widetilde u_{i,t}
+
b
\right).
\end{equation}

This cross-sectional module can capture:

\begin{itemize}
    \item common market movements;
    \item sector-level dependence;
    \item lead--lag relationships;
    \item information spillovers across firms;
    \item relative momentum and reversal;
    \item interactions between one firm's text information and related stocks.
\end{itemize}

To control complexity, the cross-sectional module will be introduced only after
a strong local stock-level model has been established.

% =========================================================
\section{iTransformer Alternative}
% =========================================================

The iTransformer reverses the conventional time-series Transformer
tokenization. Given

\begin{equation}
X
\in
\R^{L\times F},
\end{equation}

each feature trajectory is treated as one token:

\begin{equation}
x^{(f)}
=
\left[
X_{1,f},
X_{2,f},
\ldots,
X_{L,f}
\right].
\end{equation}

The trajectory is embedded as

\begin{equation}
h_f^{(0)}
=
W_{\mathrm{emb}}x^{(f)}
+
b_{\mathrm{emb}}.
\end{equation}

Self-attention is then applied across variables:

\begin{equation}
\mathcal A
=
\softmax
\left(
\frac{QK^\top}{\sqrt{d_{\mathrm{att}}}}
\right),
\qquad
\mathcal A\in\R^{F\times F}.
\end{equation}

Here $Q$ and $K$ are the projected variable-token query and key matrices,
$d_{\mathrm{att}}$ is their per-head width, and $\mathcal A$ contains the
resulting attention weights between all pairs of variables.

This explicitly models relationships such as:

\begin{equation}
\text{return}
\leftrightarrow
\text{volatility}
\leftrightarrow
\text{volume}
\leftrightarrow
\text{text-derived factors}.
\end{equation}

Two iTransformer configurations will be evaluated:

\begin{enumerate}
    \item \textbf{Within-stock iTransformer:} variables represent price,
    volatility, market, and compressed text features for one stock.

    \item \textbf{Cross-stock iTransformer:} variables represent the return
    histories of the different stocks, allowing attention to learn
    cross-sectional dependence.
\end{enumerate}

The iTransformer is trained from scratch using competition data and serves as
a locally supervised alternative to the pretrained backbone.

% =========================================================
\section{Training Objectives}
% =========================================================

\subsection{Primary Directional Loss}

The main training objective is binary cross-entropy:

\begin{equation}
\BCE
=
-
\frac{1}{B}
\sum_{(i,t)\in\mathcal{B}}
\left[
y_{i,t}\log p_{i,t}
+
(1-y_{i,t})\log(1-p_{i,t})
\right].
\end{equation}

Model selection will be based on validation hit rate, not validation
cross-entropy.

\subsection{Confidence-Weighted Classification}

Returns close to zero may have unstable labels. A weighted loss can give
greater emphasis to observations with clearer movements:

\begin{equation}
w_{i,t}
=
\min
\left(
w_{\max},
1+
\alpha
\frac{
|R_{i,t}^{(H)}|
}{
\widehat \sigma_{i,t}^{(H)}+\epsilon
}
\right).
\end{equation}

The loss becomes

\begin{equation}
\mathcal{L}_{\mathrm{weighted}}
=
-
\frac{1}{B}
\sum_{(i,t)\in\mathcal B}
w_{i,t}
\left[
y_{i,t}\log p_{i,t}
+
(1-y_{i,t})\log(1-p_{i,t})
\right].
\end{equation}

This modification will be treated as an ablation because weighting examples
can improve representation learning but may also create a mismatch with an
unweighted hit-rate metric.

\subsection{Optional Auxiliary Return Loss}

Although regression is not required for hit-rate optimization, an auxiliary
return-prediction head may regularize the latent representation:

\begin{equation}
\widehat R_{i,t}^{(H)}
=
g_R(u_{i,t}).
\end{equation}

The combined loss is

\begin{equation}
\mathcal{L}
=
\BCE
+
\lambda_R
\operatorname{Huber}
\left(
\widehat R_{i,t}^{(H)},
R_{i,t}^{(H)}
\right).
\end{equation}

The coefficient $\lambda_R$ will be selected from

\begin{equation}
\lambda_R
\in
\{0,0.01,0.05,0.1,0.2\}.
\end{equation}

The pure classification model corresponds to $\lambda_R=0$ and remains the
default.

\subsection{Distillation Objective}

When multiple allowed pretrained models are available, a student model may
learn from their directional distributions.

The student is defined as the standalone multimodal GRU introduced above. It
processes the competition-provided price and text sequence

\begin{equation}
v_{i,\tau}
=
\left[
x^P_{i,\tau};
\widetilde z^T_{i,\tau}
\right],
\qquad
\tau=t-L+1,\ldots,t,
\end{equation}

and produces

\begin{equation}
p^{\mathrm{student}}_{i,t}
=
\sigmoid\left(
w_S^\top h^{\mathrm{GRU}}_{i,t}+b_S
\right).
\end{equation}

Here $h^{\mathrm{GRU}}_{i,t}$ is the final recurrent state and $w_S$ and
$b_S$ are the student's classification-head parameters. The student does not
receive $h^{\mathrm{FM}}_{i,t}$ or any teacher output as an inference-time
input. Consequently, after training it can make predictions without running
the foundation-model teachers; otherwise, the method would be forecast
stacking rather than distillation.

Let

\begin{equation}
p^{\mathrm{teacher}}_{i,t}
=
\sum_{m=1}^{M_{\mathrm{mod}}}
\alpha_m p^{(m)}_{i,t},
\qquad
\alpha_m\geq 0,
\qquad
\sum_{m=1}^{M_{\mathrm{mod}}} \alpha_m=1.
\end{equation}

Here $M_{\mathrm{mod}}$ is the number of teacher models and $\alpha_m$ is the
nonnegative weight assigned to teacher $m$.

The student objective is

\begin{equation}
\mathcal{L}_{\mathrm{distill}}
=
\mathcal{L}^{\mathrm{student}}_{\mathrm{BCE}}
+
\frac{\lambda_D}{B}
\sum_{(i,t)\in\mathcal B}
\operatorname{KL}
\left[
\operatorname{Bern}\left(p^{\mathrm{teacher}}_{i,t}\right)
\parallel
\operatorname{Bern}\left(p^{\mathrm{student}}_{i,t}\right)
\right],
\end{equation}

where

\begin{equation}
\mathcal{L}^{\mathrm{student}}_{\mathrm{BCE}}
=
-\frac{1}{B}
\sum_{(i,t)\in\mathcal B}
\left[
y_{i,t}\log p^{\mathrm{student}}_{i,t}
+(1-y_{i,t})\log\left(1-p^{\mathrm{student}}_{i,t}\right)
\right].
\end{equation}

Thus, the first term trains the student on the realized directional labels,
whereas the second term transfers the soft probabilities of the teacher
ensemble. The operator $\operatorname{Bern}(p)$ denotes a Bernoulli
distribution with success probability $p$, and $\lambda_D\geq 0$ controls the
strength of distillation.

Distillation will only be used if the competition rules allow outputs from the
permitted models to serve as training targets.

% =========================================================
\section{Validation Protocol}
% =========================================================

Random train--validation splitting is inappropriate because four-week targets
from nearby dates overlap heavily. It would create information leakage and
overstate out-of-sample performance.

The study will therefore use purged walk-forward validation. For fold
$\kappa\in\{1,\ldots,K_{\mathrm{fold}}\}$:

\begin{align}
\mathcal D_{\mathrm{train}}^{(\kappa)}
&=
\{(i,t): t\leq T_\kappa\},\\
\mathcal D_{\mathrm{validation}}^{(\kappa)}
&=
\{(i,t): T_\kappa+G < t\leq V_\kappa\},
\end{align}

where $T_\kappa$ is the final training date, $V_\kappa$ is the final
validation date, and $G\geq H$ is the purge or embargo interval.

The following validation statistics will be reported:

\begin{itemize}
    \item overall hit rate;
    \item hit rate by stock;
    \item hit rate by sector;
    \item hit rate by calendar period;
    \item hit rate by volatility regime;
    \item balanced accuracy;
    \item upward and downward class accuracy;
    \item probability calibration;
    \item fold-to-fold standard deviation.
\end{itemize}

A robust model-selection score may be defined as

\begin{equation}
S_{\mathrm{robust}}
=
\frac{1}{K_{\mathrm{fold}}}
\sum_{\kappa=1}^{K_{\mathrm{fold}}}
\mathrm{HitRate}_\kappa
-
\gamma
\operatorname{Std}
\left(
\mathrm{HitRate}_1,
\ldots,
\mathrm{HitRate}_{K_{\mathrm{fold}}}
\right).
\end{equation}

The hyperparameter $\gamma\geq 0$ controls the penalty assigned to
fold-to-fold instability.

This penalizes models that achieve high average performance through one
exceptionally favorable validation period.

% =========================================================
\section{Threshold Selection}
% =========================================================

The default threshold $q=0.5$ may not maximize hit rate because:

\begin{itemize}
    \item equity returns may have a positive unconditional drift;
    \item predicted probabilities may be miscalibrated;
    \item class proportions may change across regimes;
    \item asymmetric model errors may shift the optimal decision threshold.
\end{itemize}

A global threshold will be selected as

\begin{equation}
q^\star
=
\arg\max_{q\in\mathcal Q}
\frac{1}{K_{\mathrm{fold}}}
\sum_{\kappa=1}^{K_{\mathrm{fold}}}
\mathrm{HitRate}_\kappa(q),
\end{equation}

where

\begin{equation}
\mathcal Q
=
\{0.35,0.36,\ldots,0.65\}.
\end{equation}

Sector-specific or regime-specific thresholds will be considered only if they
improve several validation folds consistently. Stock-specific thresholds are
likely to overfit and will not be used unless there is sufficient history.

% =========================================================
\section{Ensemble Construction}
% =========================================================

Let

\begin{equation}
p^{(1)}_{i,t},
p^{(2)}_{i,t},
\ldots,
p^{(M_{\mathrm{mod}})}_{i,t}
\end{equation}

denote out-of-fold directional predictions from different models.

A linear ensemble is

\begin{equation}
p^{\mathrm{ens}}_{i,t}
 =
\sum_{m=1}^{M_{\mathrm{mod}}}
\alpha_m p^{(m)}_{i,t},
\qquad
\alpha_m\geq 0,
\qquad
\sum_{m=1}^{M_{\mathrm{mod}}}\alpha_m=1.
\end{equation}

Each architecture--fusion combination that produces a separate out-of-fold
probability is treated as a separate ensemble member.  The candidate members
are:

\begin{itemize}
    \item LightGBM and CatBoost classifiers;
    \item TabPFN using engineered price and supplied text features;
    \item Chronos directional forecast;
    \item Moirai directional forecast;
    \item TimesFM directional forecast;
    \item TCN price encoder with concatenated text;
    \item TCN price encoder with elementwise gated text fusion;
    \item TCN price encoder with text-to-price cross-attention;
    \item frozen TimesFM plus TFT market encoder with unified raw-text
    self-attention and market-to-text cross-attention;
    \item LoRA-adapted TimesFM price encoder with concatenated text, if
    permitted;
    \item LoRA-adapted TimesFM price encoder with elementwise gated text
    fusion, if permitted;
    \item LoRA-adapted TimesFM price encoder with text-to-price
    cross-attention, if permitted;
    \item GRU-based temporal price--text fusion;
    \item train-from-scratch TFT using price and supplied text inputs;
    \item train-from-scratch PatchTST-style classifier;
    \item train-from-scratch iTransformer classifier;
    \item cross-sectional Transformer classifier.
\end{itemize}

The ensemble weights will be selected using only out-of-fold predictions.
A nonlinear stacking model may additionally use:

\begin{itemize}
    \item individual model probabilities;
    \item model disagreement;
    \item forecast dispersion;
    \item recent volatility;
    \item stock or sector indicator;
    \item strength of the supplied text signal.
\end{itemize}

Model disagreement is defined as

\begin{equation}
D_{i,t}
=
\operatorname{Std}
\left(
p^{(1)}_{i,t},
\ldots,
p^{(M_{\mathrm{mod}})}_{i,t}
\right).
\end{equation}

High disagreement may indicate low-confidence cases or regime-dependent model
performance.

% =========================================================
\section{Experimental Design}
% =========================================================

The empirical study will proceed in stages.

\subsection{Stage 1: Establish Simple Baselines}

\begin{enumerate}
    \item always-up prediction;
    \item momentum and reversal rules;
    \item logistic regression with price features only;
    \item logistic regression with price and compressed text embeddings;
    \item LightGBM using price features only;
    \item LightGBM using price and compressed text embeddings;
    \item CatBoost using price features only;
    \item CatBoost using price and compressed text embeddings;
    \item XGBoost using price features only;
    \item XGBoost using price and compressed text embeddings.
\end{enumerate}

The logistic-regression baselines will use an $L_2$ penalty, also called a
ridge or squared-coefficient penalty.  For an input vector $q_{i,t}$, where
$q_{i,t}=x^P_{i,t}$ in the price-only model and
$q_{i,t}=[x^P_{i,t};\widetilde z^T_{i,t}]$ in the multimodal model, define

\begin{equation}
p^{\mathrm{LR}}_{i,t}
=
\sigmoid\left(
\beta_0+\beta^\top q_{i,t}
\right).
\end{equation}

The coefficients are estimated by minimizing

\begin{equation}
\begin{aligned}
\mathcal L_{\mathrm{LR}}
={}&
-\frac{1}{B}
\sum_{(i,t)\in\mathcal B}
\left[
y_{i,t}\log p^{\mathrm{LR}}_{i,t}
+(1-y_{i,t})\log\left(1-p^{\mathrm{LR}}_{i,t}\right)
\right]
+\lambda_2\lVert\beta\rVert_2^2,
\end{aligned}
\end{equation}

where $\beta_0$ is the unpenalized intercept, $\beta$ is the coefficient
vector, and $\lambda_2\geq0$ controls the squared penalty.  All continuous
inputs are standardized using training-fold statistics only, and
$\lambda_2$ is selected using the purged walk-forward validation folds.  The
penalty is particularly useful for the price-plus-text model because its
compressed embedding features may remain correlated or high-dimensional.

\subsection{Stage 2: Evaluate Pretrained Forecasts}

\begin{enumerate}
    \item zero-shot Chronos;
    \item zero-shot Moirai;
    \item zero-shot TimesFM;
    \item majority vote;
    \item weighted forecast ensemble;
    \item tabular classifier using all foundation-model outputs.
\end{enumerate}

\subsection{Stage 3: Train Local Sequential Models}

\begin{enumerate}
    \item TCN with price features;
    \item TCN with concatenated text embeddings;
    \item TCN with gated text fusion;
    \item TCN with cross-attention;
    \item GRU baseline;
    \item Temporal Fusion Transformer trained from scratch;
    \item PatchTST-style classifier;
    \item iTransformer classifier.
\end{enumerate}

\subsection{Stage 4: Transfer Pretrained Representations}

\begin{enumerate}
    \item frozen TimesFM representation as static context for a trainable TFT
    over historical market covariates;
    \item family-specific adapters over original, uncompressed embeddings;
    \item joint text self-attention followed by market-to-text
    cross-attention;
    \item LoRA adaptation, if permitted;
    \item partial unfreezing of the final backbone blocks, if permitted.
\end{enumerate}

\subsection{Stage 5: Model Cross-Stock Relationships}

\begin{enumerate}
    \item cross-sectional engineered features;
    \item Transformer attention across stock representations;
    \item sector-restricted attention;
    \item cross-sectional iTransformer;
    \item local-plus-cross-sectional ensemble.
\end{enumerate}

% =========================================================
\section{Ablation Studies}
% =========================================================

The following ablations will identify which components provide incremental
directional information.

\begin{table}[ht]
\centering
\caption{Planned ablation experiments.}
\begin{tabular}{p{0.29\textwidth}p{0.61\textwidth}}
\toprule
\textbf{Component} & \textbf{Comparison} \\
\midrule
Text information
&
Price only versus price plus each supplied embedding family.
\\

Embedding fusion
&
Concatenation versus averaging versus gated fusion versus cross-attention.
\\

Pretrained information
&
No pretrained model versus forecast features versus frozen latent states.
\\

Adaptation
&
Frozen backbone versus LoRA versus partial fine-tuning.
\\

Pretrained model
&
Chronos versus Moirai versus TimesFM versus their ensemble.
\\

Temporal encoder
&
TCN versus GRU versus TFT versus PatchTST versus iTransformer.
\\

Objective
&
Binary cross-entropy versus weighted BCE versus BCE plus auxiliary regression.
\\

Context length
&
$L\in\{64,128,256,512\}$.
\\

Cross-sectional modeling
&
Independent stock predictions versus cross-stock Transformer.
\\

Threshold
&
Fixed $0.5$ versus validation-selected global threshold.
\\

Embedding dimension
&
PCA or projection dimensions in $\{16,32,64,128\}$.
\\
\bottomrule
\end{tabular}
\end{table}

% =========================================================
\section{Statistical Evaluation}
% =========================================================

For each model, the hit rate is a sample proportion. A simple standard error is

\begin{equation}
\operatorname{SE}
\left(
\widehat{\mathrm{HitRate}}
\right)
=
\sqrt{
\frac{
\widehat h(1-\widehat h)
}{
n
}
},
\end{equation}

where $\widehat h$ is the observed hit rate and $n$ is the number of predictions.

Because stock-date observations are cross-sectionally and temporally
dependent, naive independent-observation standard errors may be misleading.
The analysis will therefore also use:

\begin{itemize}
    \item date-level block bootstrap;
    \item paired comparison of daily hit rates;
    \item stock-level performance dispersion;
    \item fold-level robustness analysis.
\end{itemize}

For models $A$ and $B$, define the date-level difference

\begin{equation}
d_t
=
\frac{1}{N_t}
\sum_i
\left[
\ind\{\widehat y^A_{i,t}=y_{i,t}\}
-
\ind\{\widehat y^B_{i,t}=y_{i,t}\}
\right].
\end{equation}

Block-bootstrap confidence intervals for $\E[d_t]$ will be used to determine
whether improvements are robust rather than driven by isolated dates.

% =========================================================
\section{Expected Contributions}
% =========================================================

The proposed work offers four main contributions.

\begin{enumerate}
    \item \textbf{Metric-aligned adaptation.}
    It reformulates pretrained numerical forecasting as directional
    classification aligned with the competition's hit-rate objective.

    \item \textbf{Richer use of foundation models.}
    It compares final forecast stacking, hidden-representation extraction, and
    parameter-efficient weight adaptation rather than treating pretrained
    models only as black-box point forecasters.

    \item \textbf{Rule-compliant multimodal fusion.}
    It integrates historical market information with only the supplied
    precomputed text embeddings, without external text processing or
    prohibited language models.

    \item \textbf{Cross-sectional modeling.}
    It investigates whether relationships among the competition stocks improve
    directional prediction beyond independent per-stock forecasting.
\end{enumerate}

% =========================================================
\section{Risks and Mitigation}
% =========================================================

\subsection{Overfitting}

The primary risk is overfitting due to a limited number of distinct market
regimes. Mitigation measures include:

\begin{itemize}
    \item purged chronological validation;
    \item small trainable heads;
    \item frozen pretrained backbones;
    \item low-rank adaptation;
    \item dropout and weight decay;
    \item early stopping based on validation hit rate;
    \item avoiding stock-specific hyperparameters without strong evidence.
\end{itemize}

\subsection{Weak Foundation-Model Directional Skill}

A pretrained model may have good point-forecast MAE without useful directional
accuracy. Therefore, each foundation model will be evaluated independently on
hit rate before its outputs are included in the ensemble.

\subsection{Noisy or Redundant Embeddings}

The supplied embedding vectors may be high-dimensional, correlated, or weakly
related to four-week returns. Each embedding family will first be evaluated
separately, and low-dimensional projection or embedding dropout will be used
to reduce overfitting.

\subsection{Architectural Incompatibility}

Weights from TimesFM, Chronos, or Moirai cannot be directly copied into an
unrelated TCN, GRU, or iTransformer. Weight initialization is therefore used
only when retaining the corresponding pretrained architecture. For unrelated
student architectures, transfer will occur through forecast features,
representations, or distillation.

\subsection{Rule Interpretation}

The phrase ``use of pretrained models is allowed'' may not necessarily imply
that arbitrary fine-tuning or hidden-state extraction is allowed. The final
pipeline will include only methods explicitly compliant with the organizers'
interpretation.

% =========================================================
\section{Recommended Initial Model}
% =========================================================

The implemented high-priority model is

\begin{equation}
\boxed{
\begin{aligned}
h^P_{i,t}
&=
\operatorname{Pool}\!\left(
\operatorname{TimesFM}_{\theta_0}
(X^{\mathrm{ret20}}_{i,t-L+1:t})
\right),\\
h^C_{i,t}
&=
\operatorname{TFT}
\left([h^P_{i,t};\bar x^C_{i,t}],X^M_{i,t-31:t}\right),\\
\widehat H^T_{i,t}
&=
\operatorname{TextSelfAttention}
\left(
\left[\operatorname{Adapter}_k(e^{(k)}_{i,t,a})
+\rho_k+\eta_a\right]_{k,a}
\right),\\
h^T_{i,t}
&=
\operatorname{CrossAttention}
\left(
Q=W_Qh^C_{i,t},
K=\widehat H^T_{i,t},
V=\widehat H^T_{i,t}
\right),\\
u_{i,t}
&=
\operatorname{ResidualFusion}
\left(
[h^C_{i,t};h^T_{i,t}]
\right),\\
p_{i,t}
&=
\sigmoid
\left(
\operatorname{MLP}(u_{i,t})
\right).
\end{aligned}
}
\end{equation}

TimesFM remains frozen, whereas the TFT and the complete text/fusion path are
trained jointly. Subsequent versions may add:

\begin{enumerate}
    \item foundation-model forecast and uncertainty features;
    \item LoRA adaptation, if permitted;
    \item cross-sectional attention across stocks.
\end{enumerate}

This staged construction makes it possible to determine whether each increase
in architectural complexity produces a genuine improvement in purged
out-of-sample hit rate.

% =========================================================
\section{Implementation Plan}
% =========================================================

\begin{table}[ht]
\centering
\caption{Proposed implementation sequence.}
\begin{tabular}{clp{0.55\textwidth}}
\toprule
\textbf{Phase} & \textbf{Model} & \textbf{Purpose} \\
\midrule
1
&
Naive and tabular baselines
&
Establish unconditional drift, momentum, and text-feature benchmarks.
\\

2
&
Frozen TSFM forecasts
&
Measure zero-shot directional value of Chronos, Moirai, and TimesFM.
\\

3
&
Forecast stacking
&
Combine TSFM predictions with engineered price and text features.
\\

4
&
TCN multimodal classifier
&
Establish a compact train-from-scratch neural benchmark.
\\

5
&
Frozen TSFM representation
&
Test whether latent temporal features outperform final forecasts.
\\

6
&
Gated or attention fusion
&
Model conditional interactions between text and price patterns.
\\

7
&
LoRA adaptation
&
Adapt the pretrained backbone to the directional financial task.
\\

8
&
Cross-stock Transformer
&
Capture sector, market, and inter-stock information spillovers.
\\

9
&
Out-of-fold ensemble
&
Combine complementary model errors and improve robustness.
\\
\bottomrule
\end{tabular}
\end{table}

% =========================================================
\section{Conclusion}
% =========================================================

This proposal presents a rule-compliant framework for directional stock
forecasting using historical market variables, supplied text embeddings, and
competition-permitted time-series foundation models. The framework does not
assume that a low numerical forecasting error automatically produces a high
directional hit rate. Instead, it directly optimizes the probability of a
positive four-week return and selects all models according to purged
out-of-sample hit rate.

The main hypothesis is that pretrained time-series models are most useful not
only through their final numerical forecasts, but also through their internal
representations and carefully constrained parameter adaptation. These temporal
priors can be combined with text embeddings through gated fusion or
cross-attention and, when justified by validation, extended with attention
across stocks.

The proposed staged evaluation will determine whether the strongest practical
solution is a simple stacked classifier, a frozen-representation model, a
LoRA-adapted multimodal foundation model, or an ensemble of these approaches.


% =========================================================
\section{Methodological Novelty}
% =========================================================

The proposed work is not intended to contribute only a new architecture for
one financial forecasting competition. Its broader methodological objective is
to develop a decision-aligned and parameter-efficient adaptation framework for
pretrained time-series foundation models.

Existing time-series foundation models are generally pretrained to minimize
numerical forecasting losses. In many practical applications, however, the
forecast itself is not the final object of interest. The forecast is used to
make a downstream decision, such as determining whether a future quantity will
cross a threshold, selecting among alternatives, triggering an intervention,
or choosing an action.

This creates a potential mismatch between the pretraining objective and the
downstream evaluation criterion.

\subsection{Forecasting Loss versus Decision Loss}

Let $Y$ denote the future quantity of interest and let $\widehat Y$ denote the
model forecast. A conventional forecasting model minimizes a numerical loss

\begin{equation}
\mathcal{L}_{\mathrm{forecast}}
=
\ell(\widehat Y,Y),
\end{equation}

where $\ell$ may be the mean squared error, mean absolute error, negative
log-likelihood, or quantile loss.

The downstream decision is determined through a decision map

\begin{equation}
d:
\mathcal{Y}
\rightarrow
\mathcal{A},
\end{equation}

where $\mathcal{A}$ is the action space. In the directional forecasting task,

\begin{equation}
d(Y)
=
\ind\{Y>0\}.
\end{equation}

The corresponding decision loss is

\begin{equation}
\mathcal{L}_{\mathrm{decision}}
=
\ind
\left\{
d(\widehat Y)
\neq
d(Y)
\right\}.
\end{equation}

A model can achieve a low numerical forecasting error while producing the
wrong downstream decision. For example, when $Y$ is close to zero, even a
small forecasting error can reverse the predicted sign. Conversely, a large
magnitude error may leave the sign unchanged and therefore have no effect on
directional accuracy.

The first methodological contribution is therefore to adapt pretrained
time-series models explicitly for downstream decision quality rather than
treating numerical forecast accuracy as a sufficient surrogate.

\subsection{Decision Regret Formulation}

The framework can be generalized beyond binary direction prediction. Let

\begin{equation}
U(a,Y)
\end{equation}

denote the utility obtained from taking action $a$ when the realized outcome is
$Y$. The optimal action under perfect information is

\begin{equation}
a^\star(Y)
=
\arg\max_{a\in\mathcal A}
U(a,Y).
\end{equation}

A learned model chooses an action

\begin{equation}
a_\theta(X)
\end{equation}

from historical information $X$. The decision regret is

\begin{equation}
\mathcal{R}_{\mathrm{decision}}
\left(
a_\theta(X),Y
\right)
=
U\left(a^\star(Y),Y\right)
-
U\left(a_\theta(X),Y\right).
\end{equation}

The proposed framework minimizes expected decision regret:

\begin{equation}
\min_\theta
\E
\left[
\mathcal{R}_{\mathrm{decision}}
\left(
a_\theta(X),Y
\right)
\right].
\end{equation}

For directional classification with zero-one utility, this reduces to
minimizing sign-prediction error. In other domains, the same formulation can
represent:

\begin{itemize}
    \item whether energy demand will exceed available capacity;
    \item whether traffic volume will exceed a congestion threshold;
    \item whether inventory should be replenished;
    \item whether an anomaly alarm should be triggered;
    \item whether a physiological variable will enter a risky region;
    \item which item or asset should be ranked above another.
\end{itemize}

This formulation makes the proposed method applicable beyond financial
forecasting.

\subsection{Decision-Conditioned Parameter-Efficient Adaptation}

Let a pretrained time-series foundation model with frozen parameters
$\theta_0$ generate temporal representations

\begin{equation}
H
=
f_{\theta_0}(X)
=
\left[
h_1,
h_2,
\ldots,
h_M
\right]
\in
\R^{M\times d}.
\end{equation}

Standard low-rank adaptation modifies selected pretrained matrices using

\begin{equation}
W
=
W_0 + BA,
\end{equation}

where

\begin{equation}
A\in\R^{r\times d_{\mathrm{in}}},
\qquad
B\in\R^{d_{\mathrm{out}}\times r},
\end{equation}

and $r$ is much smaller than the original weight dimensions.

The proposed novelty is to condition the magnitude and location of the
low-rank adaptation on downstream decision relevance. Rather than applying the
same adapter uniformly to all temporal representations, the model learns which
tokens or temporal regions are most important for the decision.

Let $c$ denote contextual information, such as precomputed text embeddings,
metadata, or other permitted covariates. Define a decision-conditioning score

\begin{equation}
\alpha_j
=
\frac{
\exp
\left(
q(c)^\top k(h_j)
\right)
}{
\sum_{\ell=1}^{M}
\exp
\left(
q(c)^\top k(h_\ell)
\right)
},
\end{equation}

where $q(\cdot)$ and $k(\cdot)$ are learned projection functions.

A token-level decision-conditioned adapter is then

\begin{equation}
\Delta h_j
=
\alpha_j
B A h_j,
\end{equation}

and the adapted representation is

\begin{equation}
\widetilde h_j
=
h_j+\Delta h_j.
\end{equation}

Equivalently,

\begin{equation}
\widetilde H
=
H+
\operatorname{Diag}(\alpha)
H A^\top B^\top.
\end{equation}

This mechanism has two intended effects:

\begin{enumerate}
    \item it allocates limited adaptation capacity to temporal patterns that
    are most relevant to the downstream decision;

    \item it allows contextual information to control how the pretrained
    representation is modified, rather than only being appended after the
    frozen price encoder.
\end{enumerate}

This differs from standard cross-attention. In ordinary cross-attention,
context is used to construct a fused representation after the price
representation has already been computed. In the proposed method, context
directly governs the adaptation applied to the pretrained temporal
representation.

\subsection{Decision-Aligned Surrogate Objective}

The zero-one directional loss is non-differentiable. A smooth surrogate is
therefore required.

For a realized future return $Y$ and a scalar model score $s_\theta(X)$, define

\begin{equation}
\mathcal{L}_{\mathrm{margin}}
=
\log
\left[
1+
\exp
\left(
-
\frac{
Y\,s_\theta(X)
}{
\tau_\theta(X)+\epsilon
}
\right)
\right],
\end{equation}

where $\tau_\theta(X)>0$ is a learned local uncertainty or noise scale.

The product

\begin{equation}
Y\,s_\theta(X)
\end{equation}

is positive when the prediction and the realized outcome have the same sign.
The loss therefore encourages directional agreement. The uncertainty scale
prevents observations with intrinsically high noise from dominating the
training objective.

The full objective may combine decision quality, forecast information, and
adaptation regularization:

\begin{equation}
\mathcal{L}
=
\lambda_d
\mathcal{L}_{\mathrm{margin}}
+
\lambda_f
\mathcal{L}_{\mathrm{forecast}}
+
\lambda_c
\mathcal{L}_{\mathrm{consistency}}
+
\lambda_p
\|\phi\|_2^2,
\end{equation}

where:

\begin{itemize}
    \item $\phi$ denotes the trainable adapter and task-head parameters;
    \item $\mathcal{L}_{\mathrm{forecast}}$ preserves useful numerical
    forecasting information;
    \item $\mathcal{L}_{\mathrm{consistency}}$ prevents excessive deviation
    from the pretrained representation;
    \item $\mathcal{L}_{\mathrm{margin}}$ directly aligns the model with the
    downstream decision.
\end{itemize}

The pure decision-aligned model is obtained by setting $\lambda_f=0$.

\subsection{Representation Consistency}

Parameter-efficient adaptation may overfit the downstream dataset and destroy
transferable information learned during pretraining. To reduce catastrophic
specialization, the adapted representation can be regularized toward the
frozen representation:

\begin{equation}
\mathcal{L}_{\mathrm{consistency}}
=
\frac{1}{M}
\sum_{j=1}^{M}
\left\|
\widetilde h_j-h_j
\right\|_2^2.
\end{equation}

A relational alternative preserves pairwise token geometry:

\begin{equation}
\mathcal{L}_{\mathrm{rel}}
=
\left\|
\operatorname{sim}(\widetilde H,\widetilde H)
-
\operatorname{sim}(H,H)
\right\|_F^2,
\end{equation}

where $\operatorname{sim}$ denotes a cosine-similarity matrix.

This encourages the adapter to modify only the information needed for the
downstream decision while preserving the global structure of the pretrained
representation.

\subsection{Robust Adaptation under Temporal Distribution Shift}

Time-series applications commonly exhibit temporal distribution shift:

\begin{equation}
P_{\mathrm{train}}(X,Y)
\neq
P_{\mathrm{test}}(X,Y).
\end{equation}

The proposed framework therefore includes an optional robust decision
objective:

\begin{equation}
\min_{\phi}
\sup_{Q:
D(Q,\widehat P)\leq\rho}
\E_Q
\left[
\mathcal{L}_{\mathrm{margin}}
\left(
f_{\theta_0,\phi}(X),Y
\right)
\right],
\end{equation}

where:

\begin{itemize}
    \item $\widehat P$ is the empirical training distribution;
    \item $D$ is a statistical divergence or transport distance;
    \item $\rho$ controls the size of the ambiguity set.
\end{itemize}

A computationally simpler latent-robustness formulation is

\begin{equation}
\min_{\phi}
\E
\left[
\max_{\|\delta\|_2\leq\epsilon}
\mathcal{L}_{\mathrm{margin}}
\left(
g_\phi(H+\delta),Y
\right)
\right].
\end{equation}

The perturbation $\delta$ represents uncertainty or regime variation in the
pretrained latent space. The model is trained to preserve decision quality
under locally adverse changes in the representation.

\subsection{Theoretical Motivation}

For the directional task, let

\begin{equation}
d^\star
=
\operatorname{sign}(Y),
\qquad
\widehat d
=
\operatorname{sign}(\widehat Y).
\end{equation}

A sign error can occur only if the numerical forecast error is at least as
large as the absolute decision margin:

\begin{equation}
\ind
\left\{
\widehat d\neq d^\star
\right\}
\leq
\ind
\left\{
|\widehat Y-Y|
\geq
|Y|
\right\}.
\end{equation}

For any $a>0$,

\begin{equation}
\Prob
\left(
\widehat d\neq d^\star
\right)
\leq
\Prob(|Y|\leq a)
+
\Prob
\left(
|\widehat Y-Y|\geq a
\right).
\end{equation}

If the $p$-th moment of the forecasting error is finite, Markov's inequality
gives

\begin{equation}
\Prob
\left(
\widehat d\neq d^\star
\right)
\leq
\Prob(|Y|\leq a)
+
\frac{
\E|\widehat Y-Y|^p
}{
a^p
}.
\end{equation}

This decomposition highlights two distinct causes of directional error:

\begin{enumerate}
    \item substantial probability mass close to the decision boundary;
    \item large forecasting error away from the boundary.
\end{enumerate}

Minimizing average numerical forecasting error addresses only the second term
indirectly and may ignore the concentration of observations near the decision
boundary. The proposed margin-aware adaptation explicitly focuses model
capacity on representations that matter for downstream decisions.

A more complete theoretical analysis will investigate:

\begin{itemize}
    \item classification calibration of the proposed decision surrogate;
    \item upper bounds on decision regret;
    \item the effect of low-rank adapter complexity on generalization;
    \item robustness guarantees under bounded temporal distribution shift.
\end{itemize}

\subsection{General Research Hypotheses}

The main hypotheses are:

\begin{description}
    \item[H1.]
    Direct decision-aligned adaptation produces lower downstream decision
    regret than forecast-then-threshold approaches, even when numerical
    forecasting accuracy is similar.

    \item[H2.]
    Frozen latent representations contain more useful decision information
    than final point forecasts alone.

    \item[H3.]
    Decision-conditioned low-rank adaptation outperforms standard LoRA because
    it allocates adaptation capacity selectively across temporal tokens.

    \item[H4.]
    Conditioning adaptation on an auxiliary modality is more effective than
    introducing that modality only after temporal representation pooling.

    \item[H5.]
    Restricting adaptation through low-rank structure and representation
    consistency improves generalization under limited labels and temporal
    distribution shift.

    \item[H6.]
    The proposed method generalizes across time-series domains involving
    direction, ranking, and threshold-based decisions.
\end{description}

\subsection{Claimed Contributions}

The proposed research aims to make the following contributions:

\begin{enumerate}
    \item A general formulation of decision-aligned adaptation for pretrained
    time-series foundation models.

    \item A decision-conditioned low-rank adapter that selectively modifies
    temporal representations according to downstream decision relevance.

    \item A multimodal adaptation mechanism in which contextual embeddings
    control the adaptation of the pretrained temporal backbone.

    \item A margin-aware surrogate objective that accounts for both directional
    correctness and local outcome uncertainty.

    \item An optional robustness framework for maintaining decision quality
    under temporal distribution shift.

    \item A theoretical analysis connecting forecasting error, decision
    margins, and downstream decision regret.

    \item A multi-domain empirical evaluation across several time-series
    foundation-model backbones and downstream decision tasks.
\end{enumerate}

The intended high-level contribution can be summarized as follows:

\begin{quote}
We introduce a parameter-efficient adaptation framework that selectively
modifies pretrained time-series representations according to downstream
decision sensitivity, and we evaluate whether this approach reduces decision
regret under limited supervision and temporal distribution shift.
\end{quote}



% =========================================================
% Bibliography
% =========================================================
\bibliographystyle{plainnat}

\begin{thebibliography}{99}

\bibitem[Ansari et al.(2024)]{chronos}
Ansari, A. F., Stella, L., Turkmen, C., Zhang, X., Mercado, P.,
Shen, H., Shchur, O., Rangapuram, S. S., Arango, S. P.,
Kapoor, S., Zschiegner, J., Maddix, D. C., Wang, H.,
Mahoney, M. W., Torkkola, K., Wilson, A. G., Bohlke-Schneider, M.,
and Wang, Y.
\newblock Chronos: Learning the language of time series.
\newblock \emph{arXiv preprint arXiv:2403.07815}, 2024.

\bibitem[Das et al.(2024)]{timesfm}
Das, A., Kong, W., Sen, R., and Zhou, Y.
\newblock A decoder-only foundation model for time-series forecasting.
\newblock \emph{Proceedings of the 41st International Conference on
Machine Learning}, 2024.

\bibitem[Liu et al.(2024)]{itransformer}
Liu, Y., Hu, T., Zhang, H., Wu, H., Wang, S., Ma, L., and Long, M.
\newblock iTransformer: Inverted Transformers are effective for time series
forecasting.
\newblock \emph{International Conference on Learning Representations}, 2024.

\bibitem[Noguer i Alonso and Pereira Franklin(2026)]{financialtsfm}
Noguer i Alonso, M. and Pereira Franklin, R.
\newblock Pretrained time-series foundation models for financial return
forecasting.
\newblock \emph{arXiv preprint arXiv:2606.27100}, 2026.

\bibitem[Woo et al.(2024)]{moirai}
Woo, G., Liu, C., Kumar, A., Xiong, C., Savarese, S., and Sahoo, D.
\newblock Unified training of universal time series forecasting transformers.
\newblock \emph{Proceedings of the 41st International Conference on
Machine Learning}, 2024.

\bibitem[Bai et al.(2018)]{tcn}
Bai, S., Kolter, J. Z., and Koltun, V.
\newblock An empirical evaluation of generic convolutional and recurrent
networks for sequence modeling.
\newblock \emph{arXiv preprint arXiv:1803.01271}, 2018.

\bibitem[Cho et al.(2014)]{gru}
Cho, K., van Merri\"enboer, B., Gulcehre, C., Bahdanau, D.,
Bougares, F., Schwenk, H., and Bengio, Y.
\newblock Learning phrase representations using RNN encoder--decoder for
statistical machine translation.
\newblock \emph{Proceedings of EMNLP}, 2014.

\end{thebibliography}

\end{document}
